% !TeX root = main.tex
\label{app:conventions}
Every constant quoted in this paper depends on the following choices, which are
fixed once and used everywhere.

\paragraph{Fourier transform and the Hardy projection.}
\begin{equation}\label{eq:fourier}
 \widehat f(k)=\int_{\mathbb R}f(x)e^{-ikx}\,dx,
 \qquad
 f(x)=\frac1{2\pi}\int_{\mathbb R}\widehat f(k)e^{ikx}\,dk .
\end{equation}
With this convention the positive-energy subspace for the light-ray coordinate
used here is $\{k<0\}$, so $P_+=\mathbf1_{\{k<0\}}$ and $P_-=1-P_+$; the
distributional kernel of $P_+$ is \eqref{eq:hardy}.  (In
\cite{RigorUB} and in the Bogoliubov numerics the opposite labelling
$P_+=\mathbf1_{\{k>0\}}$ is used for the same object after $x\to-x$; the
Hilbert--Schmidt norms and determinants quoted are unaffected.)

\paragraph{Half-density push-forward.}
For an increasing $C^{1,1}$ diffeomorphism $\varphi$,
\begin{equation}\label{eq:pushforward}
 (W_\varphi f)(y)=f\bigl(\varphi^{-1}(y)\bigr)
 \bigl((\varphi^{-1})'(y)\bigr)^{1/2}
 =\frac{f(\varphi^{-1}(y))}{\sqrt{\varphi'(\varphi^{-1}(y))}} ,
\end{equation}
which is unitary on $L^2(\mathbb R)$ and isometric between the corresponding
interval subspaces.  We write $u(x)=x-\varphi(x)$ for the displacement, so
$\varphi'=1-u'$.  The recovery flow inside $D$ is
$\varphi=h_s$ of \eqref{eq:hs}, with generator $-x^2/L\cdot\partial_x$, i.e.\
$u_s(x)=sx^2/(L+sx)$ on $[0,L]$; the total first-order field is
$g_{\rm tot}=0$ on $(-a,0)$, $g_{\rm tot}(x)=-x^2/L$ on $(0,L)$, and free on
$I'=(-\infty,-a]\cup[L,\infty)$.

\paragraph{Modular frame.}
For the interval $I=(-a,L)$,
\begin{equation}\label{eq:modframe}
 \xi(x)=\log\frac{x+a}{L-x},
 \qquad
 \beta(x)=\frac{dx}{d\xi}=\frac{(x+a)(L-x)}{a+L},
 \qquad
 \beta_0=\beta(0)=\frac{aL}{a+L},
\end{equation}
and the modular modes are \eqref{eq:modes}, normalized by
$\int_{\mathbb R}d\kappa\,|\psi_\kappa\rangle\langle\psi_\kappa|=1$ on $L^2(I)$,
with $Q\psi_\kappa=q_\kappa\psi_\kappa$, $q_\kappa=(1+e^\kappa)^{-1}$; thus
$\kappa$ is the one-particle modular Hamiltonian $K=\log\frac{1-Q}Q$, and
$\kappa\to+\infty$ labels the nearly empty modes.  Modular translation is
$\xi\mapsto\xi-2\pi t$, corresponding to the KMS asymmetry
$\widehat W(-k)=e^{2\pi k}\widehat W(k)$ of the thermal spectral density at
inverse temperature $2\pi$.  A symbol $X$ with bounded kernel on $I\times I$ has
modular kernel
\begin{equation}\label{eq:modkernel}
 X(\kappa,\kappa')
 =\frac1{(2\pi)^2}\iint d\xi\,d\xi'\,
 e^{-i\kappa\xi/2\pi}\,\widetilde X(\xi,\xi')\,e^{i\kappa'\xi'/2\pi},
 \qquad
 \widetilde X(\xi,\xi')=\sqrt{\beta(x)\beta(y)}\,X(x,y),
\end{equation}
and $\Tr X=\int d\kappa\,X(\kappa,\kappa)$ when $X$ is trace class.

\paragraph{Stress tensor.}
The generator normalization is
\begin{equation}\label{eq:TT}
 \langle T(x)T(y)\rangle=\frac{c_{\rm cft}}{8\pi^2}\,(x-y-i0)^{-4},
\end{equation}
i.e.\ $T$ is normalized so that its smeared version generates the M\"obius flow
with unit coefficient.  This is the convention in which the two quadratic forms of
the recovery curve come out as
\begin{equation}\label{eq:constants}
 \boxed{\;
 g_B=\frac{2c_{\rm cft}}{3\pi^2},\qquad
 g_{\mathrm{KM}}=\frac{c_{\rm cft}}6,\qquad
 f_2=\frac{g_B}8=\frac{c_{\rm cft}}{12\pi^2},\qquad
 s_2=\frac{g_{\mathrm{KM}}}2=\frac{c_{\rm cft}}{12},\qquad
 \frac{s_2}{f_2}=\pi^2 .\;}
\end{equation}
The metric normalizations behind the factors $\frac18$ and $\frac12$ are those of
\cref{def:gB}:
$-\log\Fid=\frac{\zeta^2}8g_B+o(\zeta^2)$ and $D=\frac{\zeta^2}2g_{\mathrm{KM}}+o(\zeta^2)$ per
unit of the deformation parameter, the parameter being $\zeta$.

\paragraph{Cross-ratio variables.}
Four variables occur, all M\"obius invariants of the same configuration:
\begin{equation}\label{eq:variables}
 \eta=\frac{b(a+b+c)}{(a+b)(b+c)},\quad
 z=\frac1\eta-1=\frac{ac}{b(a+b+c)},\quad
 \eta_{\rm V}=1-\eta=\frac z{1+z},\quad
 \zeta=\frac{as}{a+L} .
\end{equation}
$\eta$ is the cross ratio of Note~1, in which the Longo--Xu conditional mutual
information is $-\frac r6\log\eta$; $\eta_{\rm V}$ is the cross ratio of
\cite{VWZ}, the complement of $\eta$; $\zeta$ is the recovery variable, with
$1+\zeta$ the cross ratio of $J_s$ inside $I$.  For the ordinary Petz map
$\zeta=2z$, for the rotated map $\zeta_t=z(1+e^{-2\pi t})$, and for the pure
geometric compression $\zeta=z$.

\paragraph{Central charge versus interval length.}
The letter $c$ denotes the length of the third interval $C$ throughout
\crefrange{sec:setting}{sec:petz}; where the central charge appears it is written
$c_{\rm cft}$, and $c_{\rm cft}=r$ for the net of $r$ complex chiral fermions.  In
\cref{prop:twobranch} the pair $(c,\bar c)$ denotes the two central charges of a
two-dimensional CFT, the interval lengths being written $L_A,L_B,L_C$ there.
